% Part: first-order-logic
% Chapter: tableaux
% Section: rules-and-proofs

\documentclass[../../../include/open-logic-section]{subfiles}

\begin{document}

\iftag{FOL}
      {\olfileid[fr]{fol}{tab}{rul}}
      {\olfileid[fr]{pl}{tab}{rul}}

\olsection{Règles et \usetoken{p}{tableau}}

!!^a{tableau} constitue une exploration systématique des différentes
manières dont !!a{sentence} peut être vraie ou fausse dans
!!a{structure}. Ses éléments de base sont des !!{signed formula}s :
des !!{sentence}s munis d'un « signe » de valeur de vérité, soit
$\True$, soit~$\False$. Ces !!{formula}s signées sont disposées dans
un arbre qui croît vers le bas.

\begin{defn}
  Une \emph{!!{signed formula}} est un couple formé d'une valeur de
  vérité et d'!!a{sentence}, autrement dit :
  \[
  \sFmla{\True}{!A} \text{ ou } \sFmla{\False}{!A}.
  \]
\end{defn}

Intuitivement, on peut lire $\sFmla{\True}{!A}$ comme « $!A$ pourrait
être vraie » et $\sFmla{\False}{!A}$ comme « $!A$ pourrait être
fausse » dans une !!{structure}.

Chaque !!{signed formula} de l'arbre est soit une \emph{hypothèse}
(les hypothèses sont placées tout en haut de l'arbre), soit obtenue
d'!!a{signed formula} située plus haut par l'une des règles d'inférence.
Deux règles correspondent à chaque !!{main operator} possible de la
!!{formula} précédente : l'une lorsque le signe est~$\True$, l'autre
lorsqu'il est~$\False$. Certaines règles ramifient l'arbre ; d'autres
se contentent d'ajouter des !!{signed formula}s à la branche.
Une règle peut, et doit souvent, être appliquée non à la !!{signed formula}
qui précède immédiatement, mais à n'importe quelle !!{signed formula}
qui figure sur la branche, entre la racine et le point d'application.

Une branche est \emph{fermée} lorsqu'elle contient à la fois
$\sFmla{\True}{!A}$ et $\sFmla{\False}{!A}$. Un !!{tableau} est fermé
lorsque toutes ses branches le sont. Selon l'interprétation intuitive,
chaque branche décrit une possibilité conjointe, mais
$\sFmla{\True}{!A}$ et $\sFmla{\False}{!A}$ y sont incompatibles.
Une branche fermée élimine donc la possibilité qu'elle
décrit. En particulier, un !!{tableau} fermé élimine toute possibilité
de rendre vraies toutes les hypothèses de la forme
$\sFmla{\True}{!A}$ et fausses toutes celles de la forme
$\sFmla{\False}{!A}$.

Un !!{tableau} fermé \emph{pour $!A$} est un !!{tableau} fermé dont la
racine est~$\sFmla{\False}{!A}$. S'il existe un tel !!{tableau}, toutes
les possibilités où $!A$ serait fausse sont éliminées ; $!A$ doit donc
être vraie dans toute !!{structure}.

\end{document}
